\documentclass[12pt]{article}
\usepackage[margin=1.5in]{geometry}
\usepackage{fontspec}
\usepackage{titling}
\setmainfont{TeX Gyre Termes}

\title{The Steps}
\author{Hylocomium splendens\\
\small stair-step moss, Andrews Forest, Oregon Cascades\\
\small Paired with Mission 1.30 --- Ranger 5}
\date{}

\begin{document}
\maketitle
\thispagestyle{empty}

\noindent I do not have a brain. I should say that first, so you understand the terms. I have no neurons, no synapses, no central processing unit. What I have is chemistry. Every cell in my body knows what to do when water arrives and what to do when water leaves. The instructions are in the proteins. The schedule is in the atmosphere. I do not think about growing. I grow when it rains.

\begin{center}
--- STEP ONE ---
\end{center}

\indent The nurse log was already here when I arrived. A Douglas-fir that fell --- I cannot tell you when, because I was not here, and I have no way to know what happened before me. The log was soft enough to hold water and firm enough to hold shape. A fragment of me --- a piece broken from a colony somewhere uphill, carried by runoff or perhaps by the foot of an elk --- landed on the upper surface of the log, in a crevice where bark had peeled away, and found moisture. I do not remember landing. I remember the first rain after landing, because that is when I began.

\begin{center}
--- STEP THIRTY-ONE ---
\end{center}

\indent This year the rain stopped early. June was dry. July was dry. August was --- I cannot describe August. The air withdrew its moisture and I withdrew my life. The fronds curled inward. The green faded to brown. The chemistry of survival engaged: sugars crystallized around my proteins, protective molecules wrapped my membranes, antioxidant enzymes positioned themselves at the gates where damage would enter when water returned.

I went dark.

\begin{center}
\textit{Then it rained.}
\end{center}

\indent The first drops hit my curled fronds and the chemistry reversed. Water entered the cells. Sugars dissolved. Proteins unfolded. Chloroplasts realigned. Within five minutes --- five minutes --- I was photosynthesizing. The power was back.

\begin{center}
--- STEP FORTY-SEVEN ---
\end{center}

\indent There is a spacecraft in orbit around the Sun. Its solar panels generate electricity --- 135 watts, enough to power everything aboard. But a short circuit means the electricity has nowhere to go. The spacecraft has been desiccated for sixty-four years. It cannot recover. There is no rain in space.

I grow another step.

\vfill
\begin{center}
\small\itshape
Ranger 5 (October 18, 1962) --- solar panel failure --- 725 km lunar miss\\
The moss recovers. The spacecraft does not.\\
Both are waiting for rain.
\end{center}

\end{document}
